% !PRINTED none

\begin{center}
\gesperrt{Fünfzehnter Abschnitt.}

\textbf{Das Formenproblem der Gruppe $G_{168}$ und die Theorie der Gleichungen siebenten Grades.}

\rule{4cm}{0.4pt}

\S.\,125.

\gesperrt{Die Resolventen des Formenproblems.}
\end{center}

Das Formenproblem für die Gruppe $G_{168}$ (\S.\,43) liefert nach der allgemeinen Theorie zunächst eine Gleichung 168\textsuperscript{sten} Grades, deren Coëfficienten Invarianten sind. Jeder Theiler der Gruppe führt aber zu einer Resolvente niedrigeren Grades, und zwar vom Grade des Index des Theilers.

Wir wollen hier nur die beiden interessantesten Fälle solcher Theiler, nämlich die Octaëdergruppe

\weq{1}{G_{24} = \chi^\lambda \omega^\mu \Theta^\nu,}

und die Gruppe 21\textsuperscript{sten} Grades

\weq{2}{G_{21} = \tau^\varrho \chi^\lambda \ (\S.\,72)}

betrachten, die uns zu Resolventen 7\textsuperscript{ten} und 8\textsuperscript{ten} Grades führen.

Was zunächst die Resolventen 7\textsuperscript{ten} Grades anlangt, so können wir bei ihrer Bildung von den Hauptaxen der Gruppe ausgehen. Eine Hauptaxe nämlich bleibt durch die Gruppe $\omega^\mu \Theta^\nu$ ungeändert, und geht durch die ganze Gruppe $G_{24}$ in drei verschiedene Linien über. Es muss also sieben Tripel von Hauptaxen geben, die durch eine Gleichung 7\textsuperscript{ten} Grades bestimmt werden.

Setzen wir nach \S.\,118, (7):

\weq{3}{\begin{aligned}
A_1 &= \alpha\gamma x_1 + \gamma\beta\varepsilon^3 x_2 + \alpha\beta\varepsilon^2 x_3 \\
A_2 &= \beta\gamma\varepsilon^3 x_1 + \alpha\beta\varepsilon^2 x_2 + \alpha\gamma x_3 \\
A_3 &= \alpha\beta\varepsilon^2 x_1 + \alpha\gamma x_2 + \beta\gamma\varepsilon^3 x_3,
\end{aligned}}
%CONT>>
